\chapter{Cálculo diferencial}\label{ch:b2:diffcalc}

El cálculo de dos variables del volumen del primer año madura hasta
convertirse en el cálculo diferencial de aplicaciones entre espacios
normados: la \emph{\omterm{def:b2:diffcalc:differential}{diferencial}} como mejor aproximación lineal, la regla
de la cadena en toda su generalidad, el teorema de simetría de Schwarz
\emph{demostrado}, las fórmulas de Taylor y el análisis completo de
segundo orden de los extremos. El teorema de la función inversa, corona
de la teoría, se enuncia con su estrategia de demostración ---un punto
fijo de Banach--- hecha explícita.

En todo el capítulo, $U$ es un subconjunto \omterm{def:b2:metric:topology}{abierto} de $\R^n$ (o de un
espacio normado; el caso de dimensión finita contiene todas las ideas)
y $f \colon U \to \R^m$.

\section{La diferencial}

\begin{definition}\label{def:b2:diffcalc:differential}
$f$ es \emph{diferenciable}\index{diferenciable} en $a$ cuando existe
una aplicación lineal (\omterm{def:b2:metric:continuity}{continua}) $\dd f_a \colon \R^n \to \R^m$
con
\[
f(a + h) = f(a) + \dd f_a(h) + o(\norm h)
\qquad (h \to 0).
\]
La aplicación $\dd f_a$, la \emph{diferencial}\index{diferencial} de
$f$ en $a$, es única; su matriz en las bases canónicas es la
\emph{matriz jacobiana}\index{matriz jacobiana} $J_f(a) =
\bigl(\frac{\partial f_i}{\partial x_j}(a)\bigr)$. La
diferenciabilidad implica la \omterm{def:b2:metric:continuity}{continuidad} y la existencia de todas las
derivadas direccionales $\dd f_a(v) = \lim_{t\to0}\frac{f(a + tv) -
f(a)}{t}$; el recíproco es falso (\cref{exo:b2:diffcalc:2}). Para $m
= 1$, $\dd f_a(h) = \langle \nabla f(a), h\rangle$: el gradiente del
primer año, entendido ahora como el vector que representa la
diferencial.
\end{definition}

\begin{theorem}[Criterio $C^1$]\label{thm:b2:diffcalc:c1}
Si todas las derivadas parciales de $f$ existen sobre $U$ y son
\omterm{def:b2:metric:continuity}{continuas} en $a$, entonces $f$ es \omterm{def:b2:diffcalc:differential}{diferenciable} en $a$. Ser ``$C^1$
sobre $U$'' ---parciales continuas--- implica, por tanto, la
\omterm{def:b2:diffcalc:differential}{diferenciabilidad} en todas partes, con \omterm{def:b2:diffcalc:differential}{diferencial} \omterm{def:b2:metric:continuity}{continua}.
\end{theorem}

\begin{proof}
Componente a componente (basta $m = 1$). La demostración del primer
año para dos variables ---moverse de una coordenada en una
coordenada, aplicar el teorema del valor medio de una variable en
cada tramo y usar la \omterm{def:b2:metric:continuity}{continuidad} de las parciales en $a$--- se
generaliza palabra por palabra a $n$ tramos:
\[
f(a + h) - f(a) = \sum_{j=1}^{n} \bigl(f(a + h^{(j)}) - f(a +
h^{(j-1)})\bigr)
= \sum_j h_j\,\frac{\partial f}{\partial x_j}(\xi_j),
\]
donde $h^{(j)}$ congela las $j$ primeras coordenadas de $h$ y $\xi_j$
está en el $j$-ésimo tramo; la \omterm{def:b2:metric:continuity}{continuidad} convierte cada
$\frac{\partial f}{\partial x_j}(\xi_j)$ en $\frac{\partial
f}{\partial x_j}(a) + o(1)$, y el error es $o(\norm h)$.
\end{proof}

\begin{theorem}[Regla de la cadena]\label{thm:b2:diffcalc:chainrule}
Si $f$ es \omterm{def:b2:diffcalc:differential}{diferenciable} en $a$ y $g$ lo es en $f(a)$, entonces $g
\circ f$ es \omterm{def:b2:diffcalc:differential}{diferenciable} en $a$ con
\[
\dd(g \circ f)_a = \dd g_{f(a)} \circ \dd f_a ,
\qquad
J_{g\circ f}(a) = J_g\bigl(f(a)\bigr)\,J_f(a) :
\]
las jacobianas se multiplican.\index{regla de la cadena}
\end{theorem}

\begin{proof}
Escríbase $f(a + h) = f(a) + \dd f_a(h) + \norm h\,\varepsilon_1(h)$
y $g(b + k) = g(b) + \dd g_b(k) + \norm k\,\varepsilon_2(k)$ con $b =
f(a)$ y $k = k(h) = \dd f_a(h) + \norm h \varepsilon_1(h)$.
Sustituyendo,
\[
g(f(a+h)) = g(b) + \dd g_b\bigl(\dd f_a(h)\bigr)
+ \norm h\,\dd g_b(\varepsilon_1(h)) + \norm{k}\,\varepsilon_2(k),
\]
y ambos términos de error son $o(\norm h)$: el primero porque $\dd
g_b$ es \omterm{def:b2:metric:continuity}{continua} y $\varepsilon_1 \to 0$; el segundo porque $\norm k
\leq C\norm h$ (aplicación lineal acotada más un término pequeño) y
$\varepsilon_2(k) \to 0$ cuando $h \to 0$.
\end{proof}

\begin{example}[Funciones radiales, de una vez por todas]\label{ex:b2:diffcalc:radial}
Sean $r(x) = \norm x_2$ sobre $\R^n\setminus\{0\}$ y $f = g \circ r$
con $g$ una función $C^1$ de una variable. Primero, $r$ es
\omterm{def:b2:diffcalc:differential}{diferenciable} fuera de $0$: de $r^2 = \sum x_i^2$,
\[
\frac{\partial r}{\partial x_i} = \frac{x_i}{r},
\qquad\text{es decir,}\qquad
\nabla r(x) = \frac{x}{\norm x} ,
\]
el vector radial \omterm{def:b2:hermitian:adjoint}{unitario} (derívese $r^2$ y divídase, o aplíquese la
regla de la cadena a $\sqrt{\cdot}$). Entonces la regla de la cadena
da, para toda función radial,
\[
\nabla f(x) = g'\bigl(\norm x\bigr)\,\frac{x}{\norm x} .
\]
Ejemplo resuelto: $g(r) = \frac1r$ da $\nabla\frac{1}{\norm x} =
-\frac{x}{\norm x^3}$, el campo del inverso del cuadrado de la
gravitación y la electrostática: de dirección radial y magnitud
$\frac{1}{\norm x^2}$. Moraleja: los gradientes de las funciones
radiales son radiales porque las curvas de nivel son esferas y el
gradiente es ortogonal a ellas; en $x = 0$, en cambio, $r$ \emph{no}
es \omterm{def:b2:diffcalc:differential}{diferenciable} (ninguna aplicación lineal candidata se ajusta a
$\norm h$ desde todas las direcciones): los perfiles radiales suaves
necesitan $g'(0) = 0$ para cruzar el origen con elegancia.
\end{example}

\begin{theorem}[Desigualdad del valor medio]\label{thm:b2:diffcalc:mvi}
Sea $f$ \omterm{def:b2:diffcalc:differential}{diferenciable} sobre $U$ y supongamos que el segmento
$\intcc{a}{b} = \{a + t(b-a)\}$ está contenido en $U$. Entonces
\[
\norm{f(b) - f(a)} \leq \norm{b - a}\,
\sup_{x \in \intcc{a}{b}} \vertiii{\dd f_x} .
\]
En particular, una aplicación \omterm{def:b2:diffcalc:differential}{diferenciable} con \omterm{def:b2:diffcalc:differential}{diferencial} nula sobre
un \omterm{def:b2:metric:topology}{abierto} \emph{\omterm{def:b2:metric:connected}{conexo}} es constante.
\end{theorem}

\begin{proof}
La función $\varphi(t) = f(a + t(b-a))$ es \omterm{def:b2:diffcalc:differential}{diferenciable} sobre
$\intcc{0}{1}$ con $\varphi'(t) = \dd f_{a + t(b-a)}(b - a)$ (regla
de la cadena), de \omterm{def:b2:nvs:norm}{norma} $\leq M\norm{b-a}$ con $M$ el supremo
mostrado. Para $f$ con valores en $\R$ concluye la desigualdad del
valor medio de una variable; para valores vectoriales, aplíquese esta
a $t \mapsto \langle u, \varphi(t)\rangle$ con $u$ el vector \omterm{def:b2:hermitian:adjoint}{unitario}
en la dirección de $f(b) - f(a)$. Constancia: es localmente constante
(segmentos dentro de bolas) más la \omterm{def:b2:metric:connected}{conexidad} (el conjunto donde $f$
vale un valor dado es \omterm{def:b2:metric:topology}{abierto} y cerrado: \cref{ch:b2:metric}).
\end{proof}

\begin{example}[Una constante de Lipschitz a partir de la desigualdad del valor medio]\label{ex:b2:diffcalc:lipschitz}
¿Es $f(x, y) = \sin x\,\sin y$ \omterm{def:b2:metric:continuity}{lipschitziana} sobre $\R^2$, y con qué
constante? Su gradiente es $\nabla f = (\cos x\sin y,\ \sin x\cos
y)$, de \omterm{def:b2:nvs:norm}{norma} al cuadrado
\[
\cos^2x\sin^2y + \sin^2x\cos^2y
\leq \sin^2 y + \cos^2y\cdot 1 = 1
\]
(acótense $\cos^2x$ y $\sin^2x$ por $1$ por separado), luego
$\vertiii{\dd f_{(x,y)}} = \norm{\nabla f} \leq 1$ en todas partes, y
el~\cref{thm:b2:diffcalc:mvi} sobre el segmento entre dos puntos
cualesquiera da
\[
\abs{f(b) - f(a)} \leq \norm{b - a}_2 :
\]
$f$ es $1$-lipschitziana, y la constante es óptima (cerca del origen,
$f(x, \tfrac\pi2) = \sin x$ tiene pendiente $1$). Moraleja: la
desigualdad del valor medio convierte una cota \omterm{def:b2:funcseq:def}{puntual} de la
\omterm{def:b2:diffcalc:differential}{diferencial} en un módulo global de \omterm{def:b2:metric:continuity}{continuidad}; la vía estándar hacia
las estimaciones \omterm{def:b2:metric:continuity}{lipschitzianas} en toda dimensión, y el motor que hay
dentro del~\cref{exo:b2:diffcalc:12}.
\end{example}

\section{Derivadas segundas}

\begin{theorem}[Schwarz]\label{thm:b2:diffcalc:schwarz}
Si $f$ es $C^2$ sobre $U$ (todas las parciales segundas existen y son
\omterm{def:b2:metric:continuity}{continuas}), entonces, para todos $i, j$:\index{teorema de Schwarz}
\[
\frac{\partial^2 f}{\partial x_i\,\partial x_j}
= \frac{\partial^2 f}{\partial x_j\,\partial x_i} .
\]
\end{theorem}

\begin{proof}
Bastan dos variables ($x = x_i$, $y = x_j$, las demás congeladas).
Consideremos la diferencia segunda
\[
\Delta(h) = f(a + h, b + h) - f(a + h, b) - f(a, b + h) + f(a,b) .
\]
Fíjese $h$ y póngase $\varphi(x) = f(x, b+h) - f(x, b)$: entonces
$\Delta(h) = \varphi(a + h) - \varphi(a)$, y dos aplicaciones del
teorema del valor medio dan
\[
\Delta(h) = h\,\varphi'(\xi)
= h\Bigl(\frac{\partial f}{\partial x}(\xi, b+h) -
\frac{\partial f}{\partial x}(\xi, b)\Bigr)
= h^2\,\frac{\partial^2 f}{\partial y\,\partial x}(\xi, \eta),
\]
con $\xi \in \intoo{a}{a+h}$ y $\eta \in \intoo{b}{b+h}$. Por
\omterm{def:b2:metric:continuity}{continuidad}, $\frac{\Delta(h)}{h^2} \to \frac{\partial^2 f}{\partial
y\partial x}(a, b)$ cuando $h \to 0$. El mismo cálculo con los
papeles de las variables intercambiados (congelando primero la
segunda variable) da $\frac{\Delta(h)}{h^2} \to \frac{\partial^2
f}{\partial x \partial y}(a,b)$: los dos límites de la misma cantidad
coinciden.
\end{proof}

\begin{example}[Por qué hace falta $C^2$: el contraejemplo de Peano]\label{ex:b2:diffcalc:peano}
Sea $f(x, y) = \dfrac{xy(x^2 - y^2)}{x^2 + y^2}$, con $f(0,0) = 0$.
Fuera del origen, $f$ es $C^\infty$; en el origen existen todas las
parciales primeras y segundas, pero las mixtas discrepan. Calculemos
a lo largo de los ejes: $f(x, 0) = f(0, y) = 0$, y para $y \neq 0$,
\[
\frac{\partial f}{\partial x}(0, y)
= \lim_{x\to0}\frac{f(x,y)}{x}
= \frac{y(0 - y^2)}{y^2} = -y ,
\qquad\text{y simétricamente}\qquad
\frac{\partial f}{\partial y}(x, 0) = x .
\]
Por tanto,
\[
\frac{\partial^2 f}{\partial y\,\partial x}(0,0)
= \frac{\dd}{\dd y}\Bigl[\frac{\partial f}{\partial
x}(0,y)\Bigr]_{y=0} = -1,
\qquad
\frac{\partial^2 f}{\partial x\,\partial y}(0,0) = +1 :
\]
las dos parciales mixtas existen y difieren. No hay contradicción con
el~\cref{thm:b2:diffcalc:schwarz}: las parciales segundas de $f$ no
son \omterm{def:b2:metric:continuity}{continuas} en $0$ (pruébese a lo largo de $y = tx$). Moraleja: el
teorema de Schwarz es un teorema genuino sobre la
\emph{\omterm{def:b2:metric:continuity}{continuidad}}, no una identidad formal, y la hipótesis ``$C^2$''
de Taylor--Young más abajo hace un trabajo real.
\end{example}

\begin{theorem}[Taylor--Young de orden 2]\label{thm:b2:diffcalc:taylor}
Sean $f \colon U \to \R$ de clase $C^2$ y $a \in U$. Entonces, cuando
$h \to 0$,
\[
f(a + h) = f(a) + \langle\nabla f(a), h\rangle
+ \frac12\, \langle H_a h,\, h\rangle + o\bigl(\norm h^2\bigr),
\]
donde $H_a = \bigl(\frac{\partial^2 f}{\partial x_i\partial
x_j}(a)\bigr)$ es la \emph{matriz hessiana}\index{matriz hessiana}
(\omterm{def:b2:quadratic:adjoint}{simétrica}, por Schwarz).
\end{theorem}

\begin{proof}
Aplíquese el teorema de Taylor--Young de una variable (volumen del
primer año) a $\varphi(t) = f(a + th)$ sobre $\intcc{0}{1}$: por la
regla de la cadena, $\varphi'(t) = \langle \nabla f(a + th),
h\rangle$ y $\varphi''(t) = \langle H_{a+th}h, h\rangle$, ambas
\omterm{def:b2:metric:continuity}{continuas} en $t$. Entonces $\varphi(1) = \varphi(0) + \varphi'(0) +
\frac12\varphi''(\theta)$ (Taylor--Lagrange) con $\theta \in
\intoo{0}{1}$, y la \omterm{def:b2:metric:continuity}{continuidad} de las parciales segundas convierte
$\varphi''(\theta)$ en $\varphi''(0) + o(1)\cdot\norm h^2$
\omterm{def:b2:funcseq:def}{uniformemente}: el desarrollo enunciado.
\end{proof}

\begin{example}[Un desarrollo por dos vías]\label{ex:b2:diffcalc:taylortwoways}
Desarrollemos $f(x, y) = \eu^x\cos y$ en el origen hasta el orden
$2$. \emph{Por composición de desarrollos de una variable:}
\[
\eu^x\cos y = \Bigl(1 + x + \frac{x^2}{2} +
o(x^2)\Bigr)\Bigl(1 - \frac{y^2}{2} + o(y^2)\Bigr)
= 1 + x + \frac{x^2 - y^2}{2} + o\bigl(\norm{(x,y)}^2\bigr) .
\]
\emph{Por derivadas parciales:} $f_x = \eu^x\cos y$, $f_y =
-\eu^x\sin y$, luego $\nabla f(0) = (1, 0)$; y $f_{xx} = f$, $f_{yy}
= -f$, $f_{xy} = -\eu^x\sin y$ dan $H_0 = \operatorname{diag}(1,
-1)$: el~\cref{thm:b2:diffcalc:taylor} reproduce $1 + x + \frac12(x^2
- y^2)$. Los dos cálculos coinciden, y la vía de la composición fue
más rápida: ninguna parcial segunda. Moraleja: el origen \emph{no} es
un punto crítico ($\nabla f \neq 0$), de modo que, pese a la hessiana
indefinida, no hay ningún punto de silla que declarar: manda el
término lineal, y el test de segundo orden solo habla en los puntos
críticos.
\end{example}

\begin{theorem}[Test de extremos de segundo orden, demostrado]\label{thm:b2:diffcalc:extrema}
Sea $f$ de clase $C^2$ cerca de un punto crítico $a$ ($\nabla f(a) =
0$), con hessiana $H = H_a$.
\begin{enumerate}
  \item Si $H$ es definida positiva, $a$ es un mínimo local estricto
        (definida negativa: máximo).
  \item Si $H$ tiene \omterm{def:b2:reduction:eigen}{valores propios} de ambos signos, $a$ es un punto
        de silla: no hay extremo.
  \item Si $H$ es singular (y semidefinida), no hay conclusión.
\end{enumerate}
El test ``$rt - s^2$'' del primer año es el caso $n = 2$: $\det H =
rt - s^2$, y el signo de la traza se lee en $r$.
\end{theorem}

\begin{proof}
Por el teorema espectral (\cref{thm:b2:quadratic:spectral}), la
\omterm{def:b2:quadratic:def}{forma cuadrática} de $H$ queda encajada entre sus \omterm{def:b2:reduction:eigen}{valores propios}
extremos: $\lambda_{\min}\norm h^2 \leq \langle Hh, h\rangle \leq
\lambda_{\max}\norm h^2$.

(1) Si $\lambda_{\min} > 0$, Taylor--Young da
\[
f(a + h) - f(a) \geq \frac{\lambda_{\min}}{2}\norm h^2 -
o(\norm h^2) > 0
\]
para $h \neq 0$ pequeño: mínimo local estricto.

(2) A lo largo de un \omterm{def:b2:reduction:eigen}{vector propio} $v_+$ con $\lambda_+ > 0$: $f(a +
tv_+) - f(a) = \frac{\lambda_+}{2}t^2 + o(t^2) > 0$ para $t$ pequeño;
a lo largo de $v_-$ con $\lambda_- < 0$, la diferencia es negativa:
ambos signos ocurren en todo entorno.

(3) $f(x,y) = x^2 + y^4$, $x^2 - y^4$ y $x^2 + y^3$ comparten la
misma hessiana singular semidefinida en $0$ con tres comportamientos
distintos.
\end{proof}

\begin{example}[Una clasificación completa, con lo global incluido]\label{ex:b2:diffcalc:quarticrun}
Clasifiquemos todos los extremos de $f(x, y) = x^4 + y^4 - 4xy$ sobre
$\R^2$. \emph{Puntos críticos:} $\nabla f = (4x^3 - 4y,\ 4y^3 - 4x) =
0$ da $y = x^3$ y $x = y^3 = x^9$, luego $x(x^8 - 1) = 0$: las
soluciones reales son $(0,0)$, $(1,1)$ y $(-1,-1)$.
\emph{Hessianas:} $H = \begin{pmatrix} 12x^2 & -4\\ -4 &
12y^2\end{pmatrix}$. En $(\pm1, \pm1)$: $\begin{pmatrix} 12 & -4\\ -4
& 12\end{pmatrix}$, de \omterm{def:b2:reduction:eigen}{valores propios} $8$ y $16$: definida positiva,
mínimos locales estrictos con $f = -2$. En $(0,0)$: $\begin{pmatrix}
0 & -4\\ -4 & 0\end{pmatrix}$, de \omterm{def:b2:reduction:eigen}{valores propios} $\pm4$: un punto de
silla. \emph{Globalidad:} de $2\abs{xy} \leq x^2 + y^2$,
\[
f(x, y) \geq x^4 + y^4 - 2(x^2 + y^2)
= (x^2 - 1)^2 + (y^2 - 1)^2 + x^2 + y^2 - 2
\xrightarrow[\norm{(x,y)}\to\infty]{} +\infty :
\]
$f$ es coerciva, así que alcanza un mínimo global (\omterm{def:b2:metric:compact}{compacidad} de los
conjuntos de subnivel), necesariamente en un punto crítico: el valor
$-2$, tanto en $(1,1)$ como en $(-1,-1)$, es el mínimo global; y no
hay máximo ($f$ no está acotada superiormente). Moraleja: el test
local clasifica candidatos, pero solo un argumento de crecimiento
convierte lo ``local'' en ``global'': el patrón en dos pasos de toda
demostración de optimización de este libro.
\end{example}

\begin{method}[Clasificar los extremos de $f \colon \R^n \to \R$]\label{met:b2:diffcalc:extrema}
\begin{enumerate}
  \item Resuélvase $\nabla f = 0$ (todos los puntos críticos; en un
        dominio con frontera, trátese la frontera aparte, como en el~\cref{exo:b2:diffcalc:7}).
  \item En cada punto crítico, calcúlense la hessiana y los signos
        de sus \omterm{def:b2:reduction:eigen}{valores propios}; en dimensión $2$ basta $\det H$ y
        $\operatorname{tr} H$: si $\det < 0$, punto de silla; si
        $\det > 0$, extremo del tipo que indique el signo de la
        traza; si $\det = 0$, el test calla y hay que estudiar $f$ a
        lo largo de curvas.
  \item Para los enunciados globales, añádase un argumento de
        \omterm{def:b2:metric:compact}{compacidad} o de coercividad ($f \to +\infty$ en el infinito, o
        un conjunto de restricciones \omterm{def:b2:metric:compact}{compacto}) y compárense después
        los valores críticos.
\end{enumerate}
\end{method}

\section{El teorema de la función inversa}

\begin{theorem}[Teorema de la función inversa]\label{thm:b2:diffcalc:inverse}
Sean $f \colon U \to \R^n$ de clase $C^1$ y $a \in U$ con $\dd f_a$
\emph{invertible}. Entonces existen entornos \omterm{def:b2:metric:topology}{abiertos} $V \ni a$ y $W
\ni f(a)$ tales que $f \colon V \to W$ es una biyección con inversa
$C^1$, y\index{teorema de la función inversa}
\[
\dd (f^{-1})_{f(x)} = (\dd f_x)^{-1} \qquad (x \in V).
\]
\end{theorem}
\admitted

\begin{remark}[Por qué es cierto: la estrategia del punto fijo]
Resolver $f(x) = y$ cerca de $a$ se reescribe como la ecuación de
punto fijo $x = x + \dd f_a^{-1}\bigl(y - f(x)\bigr) =: \Phi_y(x)$;
la aplicación $\Phi_y$ tiene \omterm{def:b2:diffcalc:differential}{diferencial} $\mathrm{id} - \dd
f_a^{-1}\dd f_x$, pequeña cerca de $a$ por la \omterm{def:b2:metric:continuity}{continuidad} de $\dd f$,
de modo que $\Phi_y$ es contractiva sobre una bola cerrada pequeña y
el teorema del punto fijo de Banach
(\cref{thm:b2:metric:banach}) proporciona la única solución local $x
= f^{-1}(y)$. La \omterm{def:b2:metric:continuity}{continuidad} y la \omterm{def:b2:diffcalc:differential}{diferenciabilidad} de la inversa se
siguen entonces de las estimaciones de la contracción. La
contabilidad completa se lleva a cabo en el tercer año; la estrategia
---y el enunciado--- se usan libremente a partir de ahora. El
\emph{teorema de la función implícita} que lo acompaña (resolver
$F(x, y) = 0$ en $y(x)$ cuando $\frac{\partial F}{\partial y}$ es
invertible) se sigue aplicando el teorema a $(x, y) \mapsto (x,
F(x,y))$.
\end{remark}

\begin{example}[Coordenadas polares]\label{ex:b2:diffcalc:polar}
$\Phi(r, \theta) = (r\cos\theta, r\sin\theta)$ tiene jacobiana
\[
J_\Phi = \begin{pmatrix} \cos\theta & -r\sin\theta\\ \sin\theta &
r\cos\theta \end{pmatrix},
\qquad \det J_\Phi = r :
\]
invertible para $r \neq 0$, de modo que $\Phi$ es un difeomorfismo
local de clase $C^1$ fuera del origen; la licencia para ``pasar a
coordenadas polares'', renovada para las integrales múltiples del~\cref{ch:b2:multint}.
\end{example}

\begin{example}[Local en todas partes, global en ninguna]\label{ex:b2:diffcalc:localnotglobal}
Sea $f(x, y) = \bigl(\eu^x\cos y,\ \eu^x\sin y\bigr)$ sobre $\R^2$.
Su jacobiana,
\[
J_f = \begin{pmatrix}
\eu^x\cos y & -\eu^x\sin y\\
\eu^x\sin y & \eu^x\cos y
\end{pmatrix},
\qquad
\det J_f = \eu^{2x} > 0 ,
\]
no se anula nunca: por el~\cref{thm:b2:diffcalc:inverse}, $f$ es un
difeomorfismo local $C^1$ en \emph{todo} punto del plano. Y sin
embargo, $f$ dista mucho de ser inyectiva: $f(x, y + 2\pi) = f(x,
y)$, así que cada valor se alcanza infinitas veces; y tampoco es
sobreyectiva, pues $\norm{f(x, y)} = \eu^x > 0$ se salta el origen.
Moraleja: el teorema de la función inversa es irreduciblemente
\emph{local}; la invertibilidad de todas las $\dd f_a$ rinde un
mosaico de inversas locales que no tiene por qué ensamblarse en una
sola. (Quien conozca los números complejos reconocerá $z \mapsto
\eu^z$; el mosaico es la familia de ramas del logaritmo.) Compárese
con el~\cref{exo:b2:diffcalc:12}, donde una hipótesis global
cuantitativa sí fuerza una única inversa global.
\end{example}

\begin{remark}[Errores frecuentes]
\emph{(i) Las derivadas direccionales son baratas; las
\omterm{def:b2:diffcalc:differential}{diferenciales}, no:} pueden existir todas las derivadas
direccionales ---e incluso no depender linealmente de la
dirección--- sin que haya \omterm{def:b2:diffcalc:differential}{diferenciabilidad}
(\cref{exo:b2:diffcalc:2}); solo el criterio $C^1$
(\cref{thm:b2:diffcalc:c1}) asciende las parciales a
\omterm{def:b2:diffcalc:differential}{diferencial}. \emph{(ii) Crítico no significa extremal:} tanto los
puntos de silla (\cref{ex:b2:diffcalc:quarticrun}) como el caso
singular mudo (\cref{thm:b2:diffcalc:extrema}~(3)) se esconden tras
$\nabla f = 0$. \emph{(iii) No hay igualdad del valor medio con
valores vectoriales:} solo sobrevive la \emph{desigualdad} del~\cref{thm:b2:diffcalc:mvi} (\cref{exo:b2:diffcalc:9}); no escribas
nunca $f(b) - f(a) = \dd f_c(b-a)$ para $f$ con valores en $\R^m$,
$m \geq 2$. \emph{(iv) La invertibilidad local no es la
inyectividad:} \cref{ex:b2:diffcalc:localnotglobal}. \emph{(v) El
gradiente pertenece al producto escalar:} $\nabla f$ es el vector que
representa $\dd f_a$ en un producto escalar elegido; cámbiese el
producto (como en el ejemplo con pesos del capítulo de \omterm{def:b2:quadratic:def}{formas
cuadráticas}) y el gradiente gira, mientras que la \omterm{def:b2:diffcalc:differential}{diferencial} ---el
objeto intrínseco--- no se mueve.
\end{remark}

\begin{remark}[Dónde se usa]
Todo lo que viene después de este capítulo es cálculo diferencial
aplicado: el capítulo de ecuaciones diferenciales lineariza flujos y
usa la fórmula del \omterm{def:b2:linalg:det}{determinante} de Liouville (demostrada en el
problema de fin de semana de este capítulo); los capítulos de curvas
y superficies estudian conjuntos de nivel y parametrizaciones a
través del teorema de la función implícita; las integrales múltiples
cambian de variable mediante jacobianas. El problema de fin de semana
desarrolla el cálculo \emph{sobre el propio espacio de matrices}
---la \omterm{def:b2:diffcalc:differential}{diferencial} del \omterm{def:b2:linalg:det}{determinante}, la de la inversa, la
\omterm{ex:b2:nvs:matrixexp}{exponencial de matrices} y el grupo ortogonal como conjunto de nivel
liso---, la sombra en el segundo año de lo que el volumen del tercer
año formaliza como variedades y grupos de Lie.
\end{remark}

\section{Ejercicios}

\begin{exercise}[$\star$]\label{exo:b2:diffcalc:1}
Calcula las \omterm{def:b2:diffcalc:differential}{matrices jacobianas} de $f(x,y) = (x^2 - y^2,\, 2xy)$ y de
$\Phi(r,\theta,z) = (r\cos\theta, r\sin\theta, z)$; ¿dónde son
invertibles las \omterm{def:b2:diffcalc:differential}{diferenciales}?
\end{exercise}

\begin{exercise}[$\star$]\label{exo:b2:diffcalc:2}
Sea $f(x,y) = \frac{x^3}{x^2 + y^2}$ (con $f(0,0) = 0$). Demuestra
que todas las derivadas direccionales de $f$ en $0$ existen, pero que
$f$ no es \omterm{def:b2:diffcalc:differential}{diferenciable} en $0$ \emph{(la aplicación $v \mapsto$
derivada direccional no es lineal)}.
\end{exercise}

\begin{exercise}[$\star$]\label{exo:b2:diffcalc:3}
Halla y clasifica los puntos críticos de $f(x, y) = x^3 + y^3 - 3xy$
usando el~\cref{thm:b2:diffcalc:extrema}, y los de $g(x,y) = x^4 +
y^4 - 2(x - y)^2$.
\end{exercise}

\begin{exercise}[$\star\star$]\label{exo:b2:diffcalc:4}
Sea $f \colon \R^n \to \R$ de clase $C^1$ y \emph{homogénea de grado
$p$}: $f(tx) = t^pf(x)$ para $t > 0$. Demuestra la identidad de Euler
\[
\langle \nabla f(x), x\rangle = p\,f(x) ,
\]
y su recíproco para funciones $C^1$ sobre $\R^n\setminus\{0\}$.
\end{exercise}

\begin{exercise}[$\star\star$]\label{exo:b2:diffcalc:5}
Sean $A$ \omterm{def:b2:quadratic:adjoint}{simétrica} y $f(x) = \frac12\langle Ax, x\rangle - \langle
b, x\rangle$. Calcula $\nabla f$ y $H_f$; ¿cuándo es $f$ convexa?
Suponiendo $A$ definida positiva, prueba que $f$ tiene un único
mínimo global en la solución de $Ax = b$: la razón de ser del
descenso de gradiente.
\end{exercise}

\begin{exercise}[$\star\star$]\label{exo:b2:diffcalc:6}
(Multiplicador de Lagrange con una restricción, demostrado a mano)
Sean $f, g$ de clase $C^1$ sobre $\R^2$ y supongamos que $f$ alcanza
en $a$ un extremo local sobre el conjunto de nivel $\{g = 0\}$, con
$\nabla g(a) \neq 0$. Demuestra que $\nabla f(a) = \lambda\nabla
g(a)$ para cierto $\lambda$. \emph{(Parametriza el conjunto de nivel
cerca de $a$ mediante el teorema de la función implícita y deriva $t
\mapsto f(\gamma(t))$.)} Aplicación: extremos de $f(x,y) = xy$ sobre
la circunferencia $x^2 + y^2 = 1$.
\end{exercise}

\begin{exercise}[$\star\star$]\label{exo:b2:diffcalc:7}
Determina los extremos de $f(x, y) = x^2 + y^2 - xy + x - y$ sobre
$\R^2$, y después su máximo y su mínimo sobre el triángulo cerrado de
vértices $(0,0)$, $(1,0)$ y $(0,1)$ \emph{(puntos críticos
interiores, luego los tres lados, luego los vértices)}.
\end{exercise}

\begin{exercise}[$\star\star\star$]\label{exo:b2:diffcalc:8}
Sea $f \colon \R^2 \to \R^2$, $f(x, y) = (x + y^2,\; y + x^2)$.
Prueba que $f$ es un difeomorfismo local cerca de $0$, calcula
$\dd(f^{-1})_{(0,0)}$ y halla el mayor $r$ tal que $\dd f$ sea
invertible sobre la bola $\norm{(x,y)}_2 < r$ \emph{(calcula $\det
J_f$)}.
\end{exercise}

\begin{exercise}[$\star\star\star$]\label{exo:b2:diffcalc:9}
(Rolle falla, el valor medio sobrevive) Da una $f \colon \R \to
\R^2$, de clase $C^1$, con $f(0) = f(2\pi)$ pero $f'(t) \neq 0$ para
todo $t$ (no hay Rolle con valores vectoriales). Verifica después
sobre tu ejemplo la \emph{desigualdad} del valor medio del~\cref{thm:b2:diffcalc:mvi}.
\end{exercise}

\begin{exercise}[$\star$]\label{exo:b2:diffcalc:10}
Calcula la \omterm{def:b2:diffcalc:differential}{diferencial} y el gradiente de $f(x) = \norm x_2^2$ y de
$g(x) = \langle Ax, x\rangle$ sobre $\R^n$ ($A$ una matriz cuadrada,
no necesariamente \omterm{def:b2:quadratic:adjoint}{simétrica}), así como la hessiana de cada una. ¿Para
qué $A$ es $g$ convexa?
\end{exercise}

\begin{exercise}[$\star\star$]\label{exo:b2:diffcalc:11}
Sea $f \colon \R^n \to \R$ de clase $C^2$. Demuestra que $f$ es
convexa si y solo si su hessiana $H_x$ es semidefinida positiva en
todo $x$ \emph{(redúcelo a una variable: $t \mapsto f(a + t(b-a))$;
usa Taylor--Lagrange en una dirección y, para el recíproco, evalúa
$\varphi''$)}.
\end{exercise}

\begin{exercise}[$\star\star\star$]\label{exo:b2:diffcalc:12}
(Un teorema de inversión global) Sea $g \colon \R^n \to \R^n$ de
clase $C^1$ con $\vertiii{\dd g_x} \leq k < 1$ para todo $x$, y sea
$f = \mathrm{id} + g$.
\begin{enumerate}
  \item Prueba que $\norm{f(x) - f(y)} \geq (1 - k)\norm{x - y}$:
        $f$ es inyectiva, con inversa \omterm{def:b2:metric:continuity}{continua} sobre su imagen.
  \item Prueba que, para cada $y \in \R^n$, la aplicación $x \mapsto
        y - g(x)$ es una contracción del espacio \omterm{def:b2:metric:complete}{completo} $\R^n$, y
        concluye por el teorema del punto fijo de Banach
        (\cref{thm:b2:metric:banach}) que $f$ es sobreyectiva.
  \item Concluye que $f$ es una biyección de $\R^n$ con inversa
        $(1-k)^{-1}$-lipschitziana: una contrapartida \emph{global}
        del~\cref{thm:b2:diffcalc:inverse} (que, por el contrario,
        es puramente local).
\end{enumerate}
\end{exercise}

\section{Problema: el cálculo de matrices --- Jacobi, exponencial y
el grupo ortogonal}

\begin{problem}\label{pb:b2:diffcalc:1}
El terreno de juego más limpio para el cálculo diferencial es el
propio espacio $\mathcal M_n(\R) \simeq \R^{n^2}$: sus aplicaciones
más naturales ---producto, inversa, \omterm{def:b2:linalg:det}{determinante}, exponencial---
tienen \omterm{def:b2:diffcalc:differential}{diferenciales} de sorprendente elegancia. Este problema las
calcula todas: la \emph{serie de Neumann}, la \omterm{def:b2:diffcalc:differential}{diferencial} de la
inversa, la \emph{fórmula de Jacobi}\index{fórmula de Jacobi} para el
\omterm{def:b2:linalg:det}{determinante} con la \emph{fórmula de Liouville} como dividendo, la
\emph{exponencial de matrices}\index{exponencial de matrices} con
$\det\eu^A = \eu^{\operatorname{tr}A}$ y, por último, el grupo
ortogonal $O_n$ como conjunto de nivel liso con las matrices
antisimétricas como espacio tangente: geometría diferencial en
embrión. En todo el problema, $\vertiii\cdot$ es la
\omterm{thm:b2:nvs:continuouslinear}{norma de operador} subordinada a $\norm\cdot_2$, y $\langle X,
Y\rangle = \operatorname{tr}(X^{\mathsf T}Y)$ es el producto escalar
de Frobenius.

\medskip
\noindent\textbf{Parte I --- La serie de Neumann.}
\begin{enumerate}
  \item Demuestra la submultiplicatividad, $\vertiii{AB} \leq
        \vertiii A\,\vertiii B$, y deduce que las aplicaciones
        polinómicas de $A$ (productos de matrices,
        \omterm{def:b2:linalg:det}{determinante}, traza) son \omterm{def:b2:metric:continuity}{continuas} sobre $\mathcal
        M_n(\R)$.
  \item Para $\vertiii X < 1$, prueba que $\sum_{k\geq0}X^k$
        converge \omterm{def:b2:series:def}{absolutamente} en $\mathcal M_n(\R)$
        (\cref{thm:b2:nvs:absoluteconvergence}), que su suma es $(I -
        X)^{-1}$ y que
        \[
        \vertiii{(I - X)^{-1}} \leq \frac{1}{1 - \vertiii X},
        \qquad
        (I - X)^{-1} = I + X + O\bigl(\vertiii X^2\bigr) .
        \]
  \item Deduce que $GL_n(\R)$ es \emph{\omterm{def:b2:metric:topology}{abierto}}: si $A$ es invertible
        y $\vertiii H < \frac{1}{\vertiii{A^{-1}}}$, entonces $A + H$
        es invertible. Deduce también que $GL_n(\R)$ es \emph{denso}
        en $\mathcal M_n(\R)$ \emph{(perturba $A$ mediante
        $\varepsilon I$: $\det(A + \varepsilon I)$ es un polinomio no
        nulo en $\varepsilon$)}.
  \item Prueba que la aplicación de inversión $\Phi(A) = A^{-1}$ es
        \omterm{def:b2:metric:continuity}{continua} sobre $GL_n(\R)$.
\end{enumerate}

\medskip
\noindent\textbf{Parte II --- Primeras \omterm{def:b2:diffcalc:differential}{diferenciales}.}
\begin{enumerate}[resume]
  \item Prueba que la aplicación $A \mapsto A^2$ es \omterm{def:b2:diffcalc:differential}{diferenciable} con
        \omterm{def:b2:diffcalc:differential}{diferencial} $H \mapsto AH + HA$, y, más en general, que $A
        \mapsto A^k$ tiene \omterm{def:b2:diffcalc:differential}{diferencial} $H \mapsto \sum_{i=0}^{k-1}
        A^iHA^{k-1-i}$. ¿Por qué \emph{no} puede escribirse $kA^{k-1}
        H$ en general?
  \item Demuestra que $\Phi(A) = A^{-1}$ es \omterm{def:b2:diffcalc:differential}{diferenciable} sobre
        $GL_n(\R)$ con
        \[
        \dd\Phi_A(H) = -A^{-1}HA^{-1}
        \]
        \emph{(escribe $(A + H)^{-1} = (I + A^{-1}H)^{-1}A^{-1}$ y
        desarrolla con la pregunta 2)}. Contrasta la fórmula con el
        caso escalar $n = 1$.
  \item Para una curva $C^1$ $t \mapsto A(t) \in GL_n(\R)$, deduce
        $\bigl(A(t)^{-1}\bigr)' = -A^{-1}A'A^{-1}$, y desarrolla $t
        \mapsto (I + tB)^{-1}$ hasta el primer orden en $t = 0$.
  \item Calcula la \omterm{def:b2:diffcalc:differential}{diferencial} de $f(A) = \operatorname{tr}(A^k)$ e
        identifica su gradiente para el producto escalar de
        Frobenius:
        \[
        \dd f_A(H) = k\operatorname{tr}\bigl(A^{k-1}H\bigr),
        \qquad
        \nabla f(A) = k\,\bigl(A^{k-1}\bigr)^{\mathsf T} .
        \]
  \item Las mismas preguntas para $f(A) = \operatorname{tr}
        (A^{\mathsf T}A) = \norm A_F^2$: \omterm{def:b2:diffcalc:differential}{diferencial}, gradiente y
        hessiana (constante); concluye que $\norm\cdot_F^2$ es
        estrictamente convexa.
\end{enumerate}

\medskip
\noindent\textbf{Parte III --- La fórmula de Jacobi.}
\begin{enumerate}[resume]
  \item Demuestra que
        \[
        \det(I + H) = 1 + \operatorname{tr}H +
        O\bigl(\vertiii H^2\bigr)
        \]
        \emph{(desarrolla $\det(e_1 + h_1, \dots, e_n + h_n)$ por
        multilinealidad en las columnas: los términos con al menos
        dos columnas $h$ son $O(\vertiii H^2)$)}: $\dd(\det)_I =
        \operatorname{tr}$.
  \item Para $A$ invertible, deduce
        \[
        \dd(\det)_A(H) = \det(A)\,
        \operatorname{tr}\bigl(A^{-1}H\bigr) .
        \]
  \item Prueba que para \emph{toda} $A$ (invertible o no),
        $\frac{\partial\det}{\partial a_{ij}}(A) = C_{ij}$, el
        cofactor $(i,j)$ \emph{(desarrollo de Laplace por la fila
        $i$)}, de modo que, con la adjunta clásica
        $\operatorname{adj}A = \operatorname{com}(A)^{\mathsf T}$:
        \[
        \dd(\det)_A(H) =
        \operatorname{tr}\bigl(\operatorname{adj}(A)\,H\bigr),
        \qquad
        \nabla(\det)(A) = \operatorname{com}(A) ,
        \]
        lo que recupera la pregunta 11 cuando $A$ es invertible
        ($\operatorname{adj}A = \det(A)A^{-1}$). Esta es la
        \emph{fórmula de Jacobi}: $\bigl(\det A(t)\bigr)' =
        \operatorname{tr}\bigl( \operatorname{adj}(A(t))\,A'(t)\bigr)$.
  \item (Fórmula de Liouville) Sea $A(t)$ una curva $C^1$ de
        matrices que satisface la ecuación diferencial lineal $A'(t)
        = M(t)A(t)$. Demuestra que
        \[
        \bigl(\det A(t)\bigr)' =
        \operatorname{tr}\bigl(M(t)\bigr)\,\det A(t),
        \qquad\text{y por tanto}\qquad
        \det A(t) = \det A(0)\,
        \exp\Bigl(\int_0^t\operatorname{tr}M\Bigr)
        \]
        \emph{(usa $\operatorname{adj}(A)\,A = \det(A)I$ y la
        invariancia \omterm{def:b2:structures:generated}{cíclica} de la traza)}: la identidad del
        wronskiano que el capítulo de ecuaciones diferenciales usará
        sin parar.
  \item Prueba que $SL_n(\R) = \{\det = 1\}$ es un conjunto de nivel
        liso: en todo $A \in SL_n(\R)$, la \omterm{def:b2:diffcalc:differential}{diferencial} $\dd(\det)_A$
        es una aplicación lineal \emph{sobreyectiva} sobre $\R$
        \emph{(evalúala en $H = \frac1nA$)}.
\end{enumerate}

\medskip
\noindent\textbf{Parte IV --- La \omterm{ex:b2:nvs:matrixexp}{exponencial de matrices}.}
\begin{enumerate}[resume]
  \item Prueba que $\eu^A = \sum_{k\geq0}\frac{A^k}{k!}$ converge
        \omterm{def:b2:series:def}{absolutamente} para toda $A$, y \omterm{def:b2:funcseq:series}{normalmente} sobre toda bola,
        con $\vertiii{\eu^A} \leq \eu^{\vertiii A}$; y que $\eu^A$
        depende continuamente de $A$.
  \item Demuestra que $AB = BA$ implica $\eu^{A+B} = \eu^A\eu^B$
        \emph{(\omterm{thm:b2:series:fubini}{producto de Cauchy}, legítimo por la convergencia
        absoluta)}; deduce que $\eu^A$ es siempre invertible, con
        inversa $\eu^{-A}$: $\exp$ manda $\mathcal M_n(\R)$ dentro de
        $GL_n(\R)$.
  \item Prueba que $t \mapsto \eu^{tA}$ es de clase $C^1$ (de hecho
        $C^\infty$) con
        \[
        \frac{\dd}{\dd t}\,\eu^{tA} = A\,\eu^{tA} =
        \eu^{tA}A
        \]
        \emph{(deriva la serie término a término sobre segmentos: la
        serie derivada converge \omterm{def:b2:funcseq:series}{normalmente})}.
  \item Demuestra la identidad
        \[
        \det\bigl(\eu^{A}\bigr) = \eu^{\operatorname{tr}A}
        \]
        \emph{(aplica la fórmula de Liouville, pregunta 13, a $A(t) =
        \eu^{tA}$)}. Comprobaciones de sensatez: $n = 1$; $A$
        nilpotente; y que las matrices de traza nula aterrizan en
        $SL_n(\R)$.
  \item Prueba que $\eu^H = I + H + O(\vertiii H^2)$, de modo que
        $\exp$ es \omterm{def:b2:diffcalc:differential}{diferenciable} en $0$ con $\dd(\exp)_0 =
        \mathrm{id}$; concluye con el teorema de la función inversa
        (\cref{thm:b2:diffcalc:inverse}) que $\exp$ es un
        difeomorfismo $C^1$ de un entorno de $0$ sobre un entorno de
        $I$: toda matriz próxima a la identidad tiene logaritmo.
  \item Prueba que $\exp$ manda las matrices \omterm{def:b2:quadratic:adjoint}{simétricas} a las
        matrices \omterm{def:b2:quadratic:adjoint}{simétricas} \emph{definidas positivas}, de manera
        biyectiva \emph{(diagonaliza; la inversa es el logaritmo
        espectral)}.
\end{enumerate}

\medskip
\noindent\textbf{Parte V --- El grupo ortogonal como conjunto de
nivel.}
\begin{enumerate}[resume]
  \item Sea $F(A) = A^{\mathsf T}A$, de $\mathcal M_n(\R)$ en las
        matrices \omterm{def:b2:quadratic:adjoint}{simétricas} $S_n$. Calcula $\dd F_A(H) = A^{\mathsf
        T}H + H^{\mathsf T}A$ y prueba que, en todo $A \in O_n =
        F^{-1}(I)$, esta \omterm{def:b2:diffcalc:differential}{diferencial} es \emph{sobreyectiva} sobre
        $S_n$ \emph{(dada $S \in S_n$, prueba con $H = \frac12 AS$)}:
        $O_n$ es un conjunto de nivel liso, de dimensión $n^2 -
        \frac{n(n+1)}2 = \frac{n(n-1)}2$.
  \item Prueba que toda curva $C^1$ $A(t) \in O_n$ con $A(0) = I$
        tiene velocidad $A'(0)$ \emph{antisimétrica}, y
        recíprocamente que, para $K$ antisimétrica, la curva
        $\eu^{tK}$ permanece en $O_n$: el espacio tangente de $O_n$
        en $I$ es exactamente el de las matrices antisimétricas.
  \item Prueba que $\det\eu^{K} = 1$ para $K$ antisimétrica
        (pregunta 18): la curva exponencial vive en el grupo de
        rotaciones $SO_n$. Calcúlala por completo para $n = 2$: con
        $J = \begin{pmatrix}0 & -1\\ 1 & 0\end{pmatrix}$, demuestra
        que
        \[
        \eu^{\theta J} = \begin{pmatrix} \cos\theta &
        -\sin\theta\\ \sin\theta & \cos\theta\end{pmatrix} :
        \]
        la \omterm{ex:b2:nvs:matrixexp}{exponencial de matrices} \emph{es} la rotación de ángulo
        $\theta$, y las definiciones por series del coseno y del seno
        reaparecen dentro de una matriz.
  \item (El truco de la densidad) Usando la densidad de $GL_n(\R)$
        (pregunta 3) y la \omterm{def:b2:metric:continuity}{continuidad}, extiende de las matrices
        invertibles a todas la identidad
        \[
        \operatorname{adj}(AB) =
        \operatorname{adj}(B)\operatorname{adj}(A)
        \]
        \emph{(para $A, B$ invertibles, ambos miembros valen
        $\det(AB)(AB)^{-1}$; y ambos son polinómicos en las
        entradas)}.
  \item Síntesis. Una frase para cada punto: (i) qué capítulos
        anteriores suministraron el motor de cada parte (la
        \omterm{def:b2:metric:complete}{completitud} y las álgebras normadas; el teorema espectral; el
        teorema de la función inversa); (ii) en qué fórmula de este
        problema se apoyará el capítulo de ecuaciones diferenciales,
        y dónde; (iii) qué dicen $\dd(\det)_I = \operatorname{tr}$ y
        $\det\eu^A = \eu^{\operatorname{tr}A}$ sobre la traza y el
        \omterm{def:b2:linalg:det}{determinante} como ``volumen infinitesimal y global''; (iv)
        qué hace el volumen del tercer año con las preguntas 21--23
        (grupos de Lie y sus álgebras de Lie).
\end{enumerate}
\end{problem}
